% Tradução brasileira revista a partir do workpass espanhol congelado;
% a fonte permanece em source_es.
% SGA 3, Exposé VIII, tradução brasileira: §3.
% Autoridade: Polo--Gille Exp8-8nov09.pdf, pp. locais 8--11,
% leitor completo pp. 624--627.

\SGARefTarget{sga3:viii:section:3}\section*{3. Propriedades de exatidão do funtor
\texorpdfstring{\(D_S\)}{D\_S}}
\addcontentsline{toc}{section}{3. Propriedades de exatidão do funtor
\texorpdfstring{\(D_S\)}{D\_S}}

\medskip
\noindent\SGARefTarget{sga3:viii:theorem:3:1}\textbf{Teorema 3.1.}
\textit{Seja \(S\) um pré-esquema, e seja}
\SGARefTarget{sga3:viii:theorem:3:1:diagram:01}
\[
  0\longrightarrow M'
   \xrightarrow{u}M
   \xrightarrow{v}M''
   \longrightarrow0
\]
\textit{uma sequência exata de grupos comutativos ordinários. Tomemos
a sequência de homomorfismos transpostos}
\SGARefTarget{sga3:viii:theorem:3:1:diagram:02}
\[
  0\longrightarrow D_S(M'')
   \xrightarrow{v^t}D_S(M)
   \xrightarrow{u^t}D_S(M')
   \longrightarrow0.
\]
\begin{enumerate}
  \item[(i)]\SGARefTarget{sga3:viii:theorem:3:1:item:i:01}
  \textit{O homomorfismo \(v^t\) induz um isomorfismo de \(D_S(M'')\)
  sobre o núcleo de \(u^t\), e \(u^t\) é fielmente plano e
  quase compacto.}

  \item[(ii)]\SGARefTarget{sga3:viii:theorem:3:1:item:ii:01}%
  \footnote{Nota do editor: o que segue foi acrescentado, e a demonstração
  foi desenvolvida em conformidade.}
  \textit{\(D_S(M')\) representa o feixe quociente
  \(D_S(M)/D_S(M'')\) para a topologia \((\mathrm{fpqc})\).}
\end{enumerate}

Denotemos por \(\mathcal M\) a família dos morfismos fielmente planos
e quase compactos. Primeiro,
\SGARefLink{sga3:viii:theorem:3:1:item:ii:01}{(ii)} deduz-se de
\SGARefLink{sga3:viii:theorem:3:1:item:i:01}{(i)} (cf. IV,
\SGARefLink{sga3:iv:corollary:4:6:5:1}{4.6.5.1}). De fato, a relação de
equivalência sobre \(D_S(M)\) definida por \(u^t\) é a mesma que a
definida pelo subgrupo
\[
  \operatorname{Ker}(u^t)=D_S(M'').
\]
Como \(u^t\in\mathcal M\), essa relação de equivalência é
\(\mathcal M\)-efetiva (cf. IV,
\SGARefLink{sga3:iv:scholium:3:3:2:1}{3.3.2.1}); portanto, \(D_S(M')\)
representa o feixe quociente para a topologia \((\mathrm{fpqc})\) (cf. IV,
\SGARefLink{sga3:iv:lemma:4:6:5}{4.6.5}).

A primeira afirmação de
\SGARefLink{sga3:viii:theorem:3:1:item:i:01}{(i)} é uma consequência
imediata da definição dos funtores \(D_S(-)\). Mais geralmente,
toda sequência exata
\SGARefTarget{sga3:viii:theorem:3:1:diagram:03}
\[
  M'\longrightarrow M\longrightarrow M''\longrightarrow0
\]
sem zero à esquerda dá origem a uma sequência transposta exata
\SGARefTarget{sga3:viii:theorem:3:1:diagram:04}
\[
  0\longrightarrow D_S(M'')
   \longrightarrow D_S(M)
   \longrightarrow D_S(M').
\]
(Isso continua válido no quadro mais geral considerado no início
do \(\SGARefLink{sga3:viii:section:1}{\S1}\).) Por outro lado, como
\(D_S(M)\) e \(D_S(M')\) são afins sobre \(S\), o morfismo \(u^t\) é
necessariamente afim e, a fortiori, quase compacto, qualquer que seja o
homomorfismo \(u:M'\to M\). A segunda afirmação de
\SGARefLink{sga3:viii:theorem:3:1:item:i:01}{(i)} deduzir-se-á, portanto,
do \SGARefLink{sga3:viii:corollary:3:2:item:a:01}{item (a)} do
\SGARefLink{sga3:viii:corollary:3:2}{corolário seguinte}.

\medskip
\noindent\SGARefTarget{sga3:viii:corollary:3:2}\textbf{Corolário 3.2.}
\textit{Seja \(S\) um pré-esquema não vazio, seja \(u:M'\to M\) um
homomorfismo de grupos comutativos ordinários e seja \(u^t:G\to G'\) o
homomorfismo transposto. Então:}
\begin{enumerate}
  \item[(a)]\SGARefTarget{sga3:viii:corollary:3:2:item:a:01}
  \textit{\(u\) é um monomorfismo se, e somente se, \(u^t\) é fielmente plano;}

  \item[(b)]\SGARefTarget{sga3:viii:corollary:3:2:item:b:01}
  \textit{\(u\) é um epimorfismo se, e somente se, \(u^t\) é um
  monomorfismo; nesse caso, \(u^t\) é, de fato, uma imersão fechada.}
\end{enumerate}

Para demonstrar \SGARefLink{sga3:viii:corollary:3:2:item:a:01}{(a)},
observe-se que, se \(u\) é um monomorfismo, então
\(\mathcal O_S(M)\), como módulo sobre \(\mathcal O_S(M')\), possui uma
base não vazia: o sistema de seções definido por qualquer sistema de
representantes de \(M\) módulo \(M'\). Por conseguinte, é fielmente plano.
Reciprocamente, se \(\mathcal O_S(M)\) é fielmente plano sobre
\(\mathcal O_S(M')\), então
\[
  u^t:\mathcal O_S(M')\longrightarrow\mathcal O_S(M)
\]
é injetivo, o que, quando \(S\ne\varnothing\), implica que
\(u:M'\to M\) é injetivo.

Para demonstrar \SGARefLink{sga3:viii:corollary:3:2:item:b:01}{(b)},
observe-se que, se \(u\) é um epimorfismo, então
\(\mathcal O_S(M')\to\mathcal O_S(M)\) é sobrejetivo, de modo que \(u^t\)
é uma imersão fechada e, a fortiori, um monomorfismo. Reciprocamente, se
\(u^t\) é um monomorfismo, então \(\operatorname{Ker}u^t\) é o grupo
unidade. Pondo \(M''=\operatorname{Coker}u\), vimos que
\[
  \operatorname{Ker}u^t\simeq D_S(M'').
\]
Por \SGARefLink{sga3:viii:proposition:2:1:item:d:01}{2.1(d)}, tem-se
\(M''=0\), e \(u\) é um epimorfismo.

O argumento habitual fornece agora a
\SGARefLink{sga3:viii:corollary:3:3}{consequência seguinte} de
\SGARefLink{sga3:viii:theorem:3:1}{3.1}.

\medskip
\noindent\SGARefTarget{sga3:viii:corollary:3:3}\textbf{Corolário 3.3.}
\textit{Seja}
\SGARefTarget{sga3:viii:corollary:3:3:diagram:00}
\[
  M'\xrightarrow{u}M\xrightarrow{v}M''
\]
\textit{uma sequência exata de grupos comutativos ordinários, e tomemos a
sequência transposta}
\SGARefTarget{sga3:viii:corollary:3:3:diagram:01}
\[
  G''\xrightarrow{v^t}G\xrightarrow{u^t}G'.
\]
\textit{Então \(v^t\) induz um morfismo fielmente plano e quase compacto
de \(G''\) para \(\operatorname{Ker}u^t\), e esse núcleo é um grupo
diagonalizável isomorfo a}
\[
  D_S(v(M))=D_S(\operatorname{Coker}u).
\]

\medskip
\noindent\SGARefTarget{sga3:viii:corollary:3:4}\textbf{Corolário 3.4.}
\textit{Seja \(S\) um pré-esquema, seja \(u:G\to H\) um homomorfismo de
pré-esquemas em grupos sobre \(S\), localmente diagonalizáveis, com \(H\) de
tipo finito sobre \(S\), e ponhamos \(G'=\operatorname{Ker}u\). Então:}
\begin{enumerate}
  \item[(a)]\SGARefTarget{sga3:viii:corollary:3:4:item:a:01}
  \textit{\(G'\) é localmente diagonalizável e, se \(G\) é de tipo
  finito sobre \(S\), \(G'\) também o é;}

  \item[(b)]\SGARefTarget{sga3:viii:corollary:3:4:item:b:01}
  \textit{o quociente \(G/G'\) «existe»: mais precisamente, a relação de
  equivalência sobre \(G\) definida por \(G'\) é \(\mathcal M\)-efetiva,
  onde \(\mathcal M\) é a família dos morfismos fielmente planos e
  quase compactos (cf. IV, \SGARefLink{sga3:iv:subsection:3:4}{3.4}). Além disso,
  \(G/G'\) é localmente diagonalizável e de tipo finito sobre \(S\);}

  \item[(c)]\SGARefTarget{sga3:viii:corollary:3:4:item:c:01}
  \textit{o homomorfismo \(u:G\to H\) fatora-se de maneira única como}
  \SGARefTarget{sga3:viii:corollary:3:4:diagram:01}
  \[
    G\xrightarrow{v}G/G'\xrightarrow{w}H,
  \]
  \textit{onde \(v\) é o homomorfismo canônico e, portanto, é
  fielmente plano e quase compacto. Além disso, \(w\) é uma imersão fechada
  e, a fortiori, um monomorfismo.}
\end{enumerate}
\textit{Por fim, o quociente}
\[
  H'=H/\operatorname{Im}w
    =\operatorname{Coker}w
    =\operatorname{Coker}u
\]
\textit{existe: mais precisamente, a relação de equivalência sobre \(H\)
definida por \(G/G'\) é \(\mathcal M\)-efetiva, e \(H'\) é de tipo
finito sobre \(S\).}

A primeira afirmação de
\SGARefLink{sga3:viii:corollary:3:4:item:c:01}{(c)} deduz-se de
\SGARefLink{sga3:viii:corollary:3:4:item:b:01}{(b)}, pela definição do
quociente \(G/G'\) (cf. IV,
\SGARefLink{sga3:iv:lemma:3:2:3}{3.2.3}).
\footnote{Nota do editor: o original foi desenvolvido no que segue.}
Mostremos que o feixe quociente \((\mathrm{fpqc})\)
\(\widetilde G/\widetilde{G'}\) é representável. Isso pode ser verificado
localmente sobre \(S\), assim como todas as demais afirmações; podemos
supor, portanto, que \(G\) e \(H\) sejam diagonalizáveis, da forma
\(D_S(M)\) e \(D_S(N)\).

Como \(H\) é de tipo finito sobre \(S\), o grupo \(N\) é de tipo finito
por \SGARefLink{sga3:viii:proposition:2:1:item:b:01}{2.1(b)}. Portanto, por
\SGARefLink{sga3:viii:corollary:1:5}{1.5}, \(u\) é definido por um
homomorfismo \(u':N\to M\). De
\SGARefLink{sga3:viii:theorem:3:1}{3.1} e
\SGARefLink{sga3:viii:corollary:3:2}{3.2} deduz-se que
\[
  G'\simeq D_S(\operatorname{Coker}u')
\]
e que \(\widetilde G/\widetilde{G'}\) é representado por
\(D_S(\operatorname{Im}u')\). Além disso, considerando a sequência exata
\SGARefTarget{sga3:viii:corollary:3:4:diagram:02}
\[
  0\longrightarrow\operatorname{Ker}u'
   \longrightarrow N
   \xrightarrow{w'}\operatorname{Im}u'
   \longrightarrow0,
\]
vê-se que \(w\) é uma imersão fechada e que
\[
  H'=H/\operatorname{Im}w=D_S(\operatorname{Ker}u').
\]
Este último é de tipo finito sobre \(S\), pois \(N\), e portanto
\(\operatorname{Ker}u'\), é de tipo finito.

\medskip
\noindent\SGARefTarget{sga3:viii:remarks:3:4}\textbf{Observações.}
O resultado da existência de quocientes de
\SGARefLink{sga3:viii:corollary:3:4}{3.4} será generalizado
consideravelmente no \(\SGARefLink{sga3:viii:section:5}{\S5}\).

Por outro lado, observe-se que, na
\SGARefLink{sga3:viii:section:3}{presente seção} e na
\SGARefLink{sga3:viii:section:2}{precedente}, a hipótese
\(S\ne\varnothing\) foi utilizada apenas para assegurar a validade de certas
recíprocas que permitem deduzir propriedades dos grupos ordinários
correspondentes a partir de certas hipóteses sobre grupos diagonalizáveis
sobre \(S\). Os resultados no sentido «direto» são válidos sem restrição
alguma sobre \(S\), e as demonstrações aqui dadas aplicam-se em geral.

\medskip
\noindent\SGARefTarget{sga3:viii:corollary:3:5}\textbf{Corolário 3.5.}
\textit{Seja \(G\) um pré-esquema em grupos diagonalizável sobre \(S\), e seja
\(n\ne0\) um inteiro. Então o subgrupo \({}_nG\) de \(G\), núcleo do
homomorfismo}
\[
  n\mathbin{\cdot}\operatorname{id}_G:G\longrightarrow G,
\]
\textit{é inteiro sobre \(S\), e é finito sobre \(S\) se \(G\) é de
tipo finito sobre \(S\).}

De fato, se \(G=D_S(M)\), então
\[
  {}_nG=D_S(M/nM)
\]
por \SGARefLink{sga3:viii:theorem:3:1}{3.1}, e o resultado deduz-se de
\SGARefLink{sga3:viii:proposition:2:1:item:b:01}{2.1(b)},
\SGARefLink{sga3:viii:proposition:2:1:item:c:01}{(c)} e
\SGARefLink{sga3:viii:proposition:2:1:item:cprime:01}{\((c')\)}.
